\section{Natural coordinates for physical interval operators}

The following identities give the coordinate and placement steps supporting
\cite[Proposition~5.7]{Arad2016RigorousRG}. All endpoints and dimensions below
are nonnegative integers. Intervals are half-open and intersected with the
finite chain. Subtraction of natural numbers is truncated at zero.

\begin{definition}[Increasing coordinates on an interval]
    \label{def:physical_interval_increasing_coordinates}
    \lean{TNLean.MPS.AreaLaw.intervalShiftOrderIso}
    For $a\le b\le T$, the increasing coordinate identification is the
    order isomorphism
    \begin{align}
        e_{a,b}:\{0,\ldots,b-a-1\} & \longrightarrow[a,b), &
        e_{a,b}(i)                 & =a+i.
        \label{eq:physical_interval_increasing_coordinates}
    \end{align}
    Its inverse sends a site $x$ to $x-a$. Empty intervals are included.
\end{definition}

\begin{theorem}[A subinterval in increasing coordinates]
    \label{thm:physical_interval_subinterval_coordinates}
    \lean{TNLean.MPS.AreaLaw.intervalShiftOrderIso_mem_subinterval}
    \uses{def:physical_interval_increasing_coordinates}
    Suppose $c\le d\le T$ and $c\le u$. For each $i<d-c$,
    \begin{align}
        e_{c,d}(i)\in[u,v)
        \quad\Longleftrightarrow\quad
        u-c\le i<(u-c)+(v-u).
        \label{eq:physical_interval_subinterval_coordinates}
    \end{align}
    No ordering of $u,v$ is required: if $v<u$, both conditions are false.
\end{theorem}
\begin{proof}
    Substitute $e_{c,d}(i)=c+i$ and compare the endpoint inequalities.
\end{proof}

\begin{theorem}[The length of the natural subinterval]
    \label{thm:physical_interval_subinterval_length}
    \lean{TNLean.MPS.AreaLaw.intervalShift_subinterval_length_le}
    If $c\le u\le d$ and $v\le d$, then
    \begin{align}
        (u-c)+(v-u)\le d-c.
        \label{eq:physical_interval_subinterval_length}
    \end{align}
    This includes $v<u$, when the subinterval length is zero.
\end{theorem}
\begin{proof}
    If $u\le v$, the left side equals $v-c$. If $v<u$, it equals $u-c$.
    Each is at most $d-c$.
\end{proof}

\begin{definition}[Coordinates on the single retained complement]
    \label{def:physical_interval_complement_coordinates}
    \lean{TNLean.MPS.AreaLaw.intervalComplementEquiv}
    \uses{def:physical_interval_increasing_coordinates,
        thm:physical_interval_subinterval_coordinates}
    For $c\le d\le T$ and $c\le u$, restriction of $e_{c,d}$ gives a
    bijection
    \begin{align}
        \{i<d-c:i\notin[u-c,(u-c)+(v-u))\}
        \longrightarrow
        [c,d)\setminus
        [u,v).
        \label{eq:physical_interval_complement_coordinates}
    \end{align}
    The retained sites form one complement; a site belonging to overlapping
    endpoint regions occurs only once.
\end{definition}

\begin{definition}[Identity extension of the supplied core tensors]
    \label{def:physical_interval_identity_extension}
    \lean{MPOTensor.identityExtendInterval}
    Let $A_r(i,j)$, $0\le r<L$, be the supplied $B\times B$ site matrices,
    and suppose $a+L\le T$. Their extension to the $T$-site chain is
    \begin{align}
        \widehat A_v(i,j)=
        \begin{cases}
            A_{v-a}(i,j),   & a\le v<a+L,     \\
            \delta_{ij}I_B, & v\notin[a,a+L).
        \end{cases}
        \label{eq:physical_interval_identity_extension}
    \end{align}
    The virtual space remains the supplied $B$-dimensional space.
\end{definition}

\begin{theorem}[Reading a nested physical interval in natural coordinates]
    \label{thm:physical_interval_nested_trace}
    \lean{MPOTensor.trace_identityExtendInterval_nested_natural}
    \uses{def:physical_interval_increasing_coordinates,
        def:physical_interval_identity_extension}
    Suppose $c\le a\le b\le d\le T$. Let $A_r(i,j)$ be the supplied core
    tensors on $b-a$ sites, with scalar closing matrix $C$ on the same
    $B$-dimensional virtual space. Let $H_K$ be their trace-form operator
    placed on $K=[a,b)$ using $e_{a,b}$, and extend it canonically to
    $E=[c,d)$. Reading this operator in the increasing coordinates of $E$
    gives the matrix with entries
    \begin{align}
        \operatorname{tr}\left(C\prod_{v=0}^{d-c-1}
        \widehat A_v(\sigma_v,\tau_v)\right),
        \label{eq:physical_interval_nested_trace}
    \end{align}
    where the identity extension starts at $a-c$ inside the natural chain
    of length $d-c$. The closing matrix and virtual space are unchanged.
    Empty intervals and zero physical or virtual dimension are included.
\end{theorem}
\begin{proof}
    When the two physical configurations agree outside the core interval,
    all exterior letters are identities, and the ordered product is the
    original core product. Otherwise one exterior letter is zero. This
    is exactly the canonical extension of the core operator. Composition
    of the increasing interval inclusions gives the stated physical
    placement, and reading along the inverse increasing coordinates
    recovers its natural matrix.
\end{proof}

\begin{theorem}[A physical complement operator as a retained matrix]
    \label{thm:physical_interval_retained_matrix}
    \lean{TNLean.MPS.AreaLaw.exists_retainedMatrix_of_physical_localLift}
    Let $V$ be a finite physical site set and $e:\{0,\ldots,T-1\}\to V$
    a bijection. Fix $a,L$ with $a+L\le T$, and suppose the physical set
    $C_{\rm phys}$ consists exactly of the images under $e$ of the sites
    outside $[a,a+L)$. Let $K_{\rm phys}$ be any operator on its
    configuration space, with local dimension $q$.
    There is a matrix $K_{\rm ret}$ on configurations of the single
    natural complement such that the canonical extension of $K_{\rm phys}$,
    first read along $e$ and then split into retained and interior
    configurations, is
    \begin{align}
        K_{\rm ret}\otimes I_{q^L}.
        \label{eq:physical_interval_retained_matrix}
    \end{align}
    The retained matrix is obtained from the supplied regional matrix
    by this coordinate identification. No interior operator is required.
\end{theorem}
\begin{proof}
    \uses{def:retained_interval_coordinates}
    Relabel the canonical extension along $e$ using covariance of the
    regional lift. The resulting natural complement matrix is reindexed
    by the complement configuration bijection. Splitting a full
    configuration into its complement and interval parts leaves the
    complement entry unchanged and contributes an identity entry on the
    interval.
\end{proof}

\begin{theorem}[A trace-chain representation on a supported interval]
    \label{thm:supported_interval_trace_chain}
    \lean{TNLean.MPS.AreaLaw.exists_traceChainOperator_of_supported_interval}
    \uses{def:chain_model,def:trace_chain_operator,
        def:physical_interval_increasing_coordinates}
    Let $0\le c\le d\le L$ and let $A$ be supported on the interval
    $J=\{c,\ldots,d-1\}$. There are a finite bond dimension, site tensors
    on a chain of $d-c$ sites and a closing matrix whose trace contraction,
    transported by $i\mapsto c+i$ and extended by the identity outside $J$,
    equals $A$. Empty intervals and zero local dimension are included.
\end{theorem}
\begin{proof}
    \uses{thm:finite_matrix_trace_chain,def:physical_interval_increasing_coordinates}
    The support hypothesis supplies a matrix on the interval coordinates.
    Pull it back along the increasing identification, represent that finite
    matrix by a trace chain, and transport the representation to its sites.
    This supplies no efficient bond bound.
\end{proof}

\begin{theorem}[Inward effective cuts with a bounded collar]
    \label{thm:inward_effective_cuts}
    \lean{TNLean.MPS.AreaLaw.inward_effective_cuts}
    \uses{def:chain_model}
    Let $\ell\ge1$, $0\le p\le r\le L$ and $r-p>4\ell$.
    Put $p'=\max(p,\ell)$, $r'=\min(r,L-\ell)$,
    $a=p'-\ell$ and $b=r'-\ell-1$. Then
    \begin{align}
        a+\ell  & =p',               & b+\ell+1 & =r',\notag        \\
        a+2\ell & \le b,             & b+2\ell  & \le L-1,\notag    \\
        p       & \le p'\le r'\le r, & r'-p'    & \ge2\ell+1,\notag \\
        p'-p    & \le\ell,           & r-r'     & \le\ell.
    \end{align}
    Also $(p'-p)+(r-r')\le2\ell$ and, for
    $J_0=\{p',\ldots,r'-1\}$ and $J=\{p,\ldots,r-1\}$,
    $J_0\subseteq J$ and $|J\setminus J_0|\le2\ell$.
\end{theorem}
\begin{proof}
    The long-interval assumption allows moving each endpoint inward by at
    most $\ell$ while retaining length at least $2\ell+1$.
    Substitute the definitions in the endpoint inequalities. The removed
    collar has $(p'-p)+(r-r')$ sites, giving its cardinal bound.
\end{proof}
